% Propietario conservado: /Users/ruben/Documents/New project/output/EXTREMA_MEDIA_RAZON_RECIPROCIDAD_ES_EN_20260918/ARTICULO/ES/source/sections/nucleo.tex
\section{Núcleo formal de la Holografía Modular Triádica}
\label{x_reciprocidad:sec:nucleo}

La Holografía Modular Triádica, abreviada HMT, comienza con una estructura finita de posiciones y operaciones. El orden de construcción importa: las posiciones reciben evaluaciones aritméticas, las transiciones adquieren una firma ternaria y la dinámica conserva la información necesaria para prolongarlas. Sólo después se forman las realizaciones numéricas. Una constante reconocida al final de este proceso no sustituye a la regla que produjo su región ni a la historia que determina su evaluación.

La denominación \emph{All Possible Paths}, APP, expresa el papel de los caminos sobre el soporte. Se adopta como denominación de la construcción, con una referencia conceptual a la formulación por caminos de la mecánica cuántica; no supone que su ley aritmética sea la integral de caminos de Feynman. El \emph{Topological Phase Kernel}, TPK,\footnote{Los autores dedican asimismo las iniciales TPK a Tan Pengkiang. Esta dedicación es independiente de la definición matemática.} designa la composición de selección, transporte, actualización de memoria y prolongación. Las siglas se mantienen únicamente después de haber definido los objetos que representan.

\subsection{APP: soporte posicional y evaluaciones aritméticas}

Consideremos las nueve marcas interiores de una recta graduada entre cero y diez,
\[
 D_9=\{1,2,3,4,5,6,7,8,9\}.
\]
Su disposición horizontal y vertical produce las ochenta y una posiciones de \(D_9^2\). Las marcas interiores no deben confundirse con los diez intervalos unitarios de la recta de referencia. Esta disposición fija un alfabeto y dos direcciones; todavía no utiliza una expansión real para decidir una trayectoria.

La identificación cíclica de las posiciones conduce al soporte
\[
 X_9=(\ZZ/9\ZZ)^2.
\]
La carta de representantes positivos conserva el símbolo nueve para la clase nula. Para todo entero positivo se definen
\begin{equation}
 \rho_9(n)=1+((n-1)\bmod9),\qquad
 q_9(n)=\frac{n-\rho_9(n)}9,
 \qquad n=\rho_9(n)+9q_9(n).
 \label{x_reciprocidad:eq:residuo-cociente}
\end{equation}
La raíz digital positiva es, por tanto, una elección de representante de la clase residual. No conserva por sí sola el entero anterior a la reducción; la coordenada \(q_9\) proporciona precisamente la información que falta.

En cada posición \((i,j)\), la acumulación de \(i\) repetida \(j\) veces produce \(ij\). Las dos evaluaciones enteras y sus proyecciones son
\[
 S(i,j)=i+j,\qquad P(i,j)=ij,\qquad
 \Sigma(i,j)=\rho_9(i+j),\quad \Pi(i,j)=\rho_9(ij).
\]
La simetría \((i,j)\mapsto(j,i)\) intercambia las direcciones de acumulación y conserva ambas evaluaciones. Es una propiedad de un único soporte, no una identificación posterior entre dos matrices independientes.

\begin{definition}[Evaluación aritmética levantada]
La evaluación de APP que conserva las dos hojas es
\begin{equation}
 \mathcal A(i,j)=
 \bigl((R^+_{ij},q^+_{ij}),(R^\times_{ij},q^\times_{ij})\bigr)
 =\bigl((\rho_9(i+j),q_9(i+j)),(\rho_9(ij),q_9(ij))\bigr).
 \label{x_reciprocidad:eq:app-levantada}
\end{equation}
El término \emph{hoja} indica aquí una evaluación sobre el mismo soporte: aditiva o multiplicativa. Las dos coordenadas de cociente pertenecen al estado aritmético y no se descartan al reducir los residuos.
\end{definition}


% BEGIN OWNER /Users/ruben/Documents/New project/output/EXTREMA_MEDIA_RAZON_RECIPROCIDAD_ES_EN_20260918/ARTICULO/ES/source/figures/app.tex
\begin{figure}[tbp]
\centering
\begin{tikzpicture}[x=.53cm,y=.53cm,font=\small]
\fill[black!5] (1.5,-.5) rectangle (2.5,8.5);
\fill[black!5] (4.5,-.5) rectangle (5.5,8.5);
\fill[black!5] (-.5,2.5) rectangle (8.5,3.5);
\fill[black!5] (-.5,5.5) rectangle (8.5,6.5);
\fill[black!12] (7.5,-.5) rectangle (8.5,8.5);
\fill[black!12] (-.5,-.5) rectangle (8.5,.5);
\foreach \i in {1,...,9}{
 \node at (\i-1,9.05) {\(\i\)};
 \node at (-1.05,9-\i) {\(\i\)};
 \foreach \j in {1,...,9}{
  \pgfmathtruncatemacro{\v}{mod(\i*\j-1,9)+1}
  \node at (\j-1,9-\i) {\(\v\)};
 }
}
\foreach \k in {0,...,9}{
 \draw[black!25] (\k-.5,-.5)--(\k-.5,8.5);
 \draw[black!25] (-.5,\k-.5)--(8.5,\k-.5);
}
\draw[line width=.65pt] (2.5,3.5) rectangle (4.5,5.5);
\node[anchor=west,align=left,text width=6.5cm] at (10,7.2) {Un soporte: \(81\) posiciones.\\[5pt]Dos evaluaciones:\\ \(S(i,j)=i+j\), \(P(i,j)=ij\).\\[5pt]Cada fila se genera por acumulación:\\ \(P(i,j+1)-P(i,j)=i\).};
\node[anchor=west,align=left,text width=6.5cm] at (10,1.8) {La carta representa \(0\pmod9\) por \(9\).\\[5pt]El bloque señalado es\\ \(B_c=\left(\begin{smallmatrix}7&2\\2&7\end{smallmatrix}\right)\),\\ con autovalores \(9\) y \(5\).};
\end{tikzpicture}
\caption{Tabla multiplicativa de APP en representantes positivos. Las filas y columnas no unitarias se distinguen por trama gris; el marco identifica el bloque diagonal adyacente de diagonales iguales. La tabla dibujada es un observable del toro aritmético; sus bordes de representación no constituyen una frontera topológica intrínseca del toro.}
\label{x_reciprocidad:fig:app}
\end{figure}

% END OWNER


\begin{proposition}[Reconstrucción de las evaluaciones]
Los dos pares de \eqref{x_reciprocidad:eq:app-levantada} determinan exactamente \(i+j\) e \(ij\). Los residuos aislados no determinan esas evaluaciones enteras ni sus cocientes.
\end{proposition}
\begin{proof}
La primera afirmación es \eqref{x_reciprocidad:eq:residuo-cociente} aplicada a cada hoja. Para la segunda, \(10=1+9\) y \(19=1+2\cdot9\) tienen el mismo residuo y distinto cociente. En el soporte bidimensional, las posiciones \((5,5)\) y \((8,2)\) dan la misma suma y el mismo residuo producto, pero productos \(25=7+18\) y \(16=7+9\). La proyección residual identifica datos que la evaluación levantada distingue.
\end{proof}

\subsubsection{Grafo de Cayley, topología toroidal y grado de enrollamiento}

Las traslaciones cardinales están dadas, en coordenadas de fila y columna, por
\[
 \nu_N=(-1,0),\quad\nu_E=(0,1),\quad
 \nu_S=(1,0),\quad\nu_O=(0,-1),\qquad T_d(x)=x+\nu_d.
\]
La primera coordenada aumenta hacia el sur y la segunda hacia el este. El grafo
\[
 \mathscr G_9=\operatorname{Cay}(X_9,\{\nu_N,\nu_E,\nu_S,\nu_O\})
\]
es la retícula cíclica bidimensional: tiene ochenta y un vértices y cuatro aristas orientadas salientes por vértice. Es un grafo de Cayley para la estructura \emph{aditiva} del soporte. La multiplicación de residuos tiene divisores de cero y no convierte todas las filas de APP en un grupo multiplicativo.

Una palabra cardinal \(w=d_1\cdots d_r\) y un origen \(x_0\) determinan
\[
 x_k=x_0+\sum_{j=1}^k\nu_{d_j}\pmod9.
\]
Existen \(81\cdot4^r\) caminos orientados de longitud \(r\), incluidos retornos y retrocesos. Cuando el camino se cierra, su desplazamiento entero previo a la reducción define
\begin{equation}
 \operatorname{wind}(w)=\frac19
 \bigl(\#S(w)-\#N(w),\#E(w)-\#O(w)\bigr)\in\ZZ^2.
\end{equation}
La inversión del camino cambia el signo de este vector. El retorno a la misma posición residual no obliga, por tanto, a que el enrollamiento sea nulo. Éste es el primer ejemplo de una diferencia que reaparecerá en el retorno de fase: una coordenada puede cerrarse mientras otra registra una evolución no trivial.

\subsubsection{Cociente de transporte e identidad de cociclo}

Sea \(s_9:\ZZ/9\ZZ\to D_9\) la sección de representantes positivos. El cociente de transporte asociado a la composición aditiva, denominado también \emph{acarreo} en aritmética posicional, es
\[
 c_9(a,b)=\frac{s_9(a)+s_9(b)-s_9(a+b)}9.
\]
El carácter composicional de esta memoria queda expresado por una identidad exacta.

\begin{lemma}[Identidad de cociclo]
Para \(a,b,c\in\ZZ/9\ZZ\),
\[
 c_9(a,b)+c_9(a+b,c)=c_9(b,c)+c_9(a,b+c).
\]
\end{lemma}
\begin{proof}
Ambas sumas son iguales a
\[
 \frac{s_9(a)+s_9(b)+s_9(c)-s_9(a+b+c)}9.
\]
La asociatividad de la suma de representantes levantados requiere precisamente este término adicional al componer en el cociente.
\end{proof}

Con representantes positivos, este cociclo no está normalizado en la identidad: \(c_9(0,a)=1\). Si \(s_0\) toma representantes en \(\{0,\ldots,8\}\), su acarreo normalizado \(c_0\) satisface
\[
 c_9(a,b)=c_0(a,b)+\mathbf1_{a=0}+\mathbf1_{b=0}-\mathbf1_{a+b=0}.
\]
La diferencia es el cambio de sección; no altera la identidad de cociclo. La descomposición posterior del calendario en fase y ventana utilizará los representantes de cero a ocho.

Los totales de ambas evaluaciones distinguen la contribución visible y la conservada en el cociente:
\[
 \sum S=810,\quad\sum P=2025,\quad
 \sum\Sigma=405,\quad\sum\Pi=459,
\]
\[
 \sum q^+=45,\qquad\sum q^\times=174.
\]
De aquí se obtiene
\begin{equation}
 2025-810=(459-405)+9(174-45)=54+9\cdot129.
 \label{x_reciprocidad:eq:defecto-integrado}
\end{equation}
La diferencia suma--producto no se agota en el defecto visible \(54\). La identidad conserva también el defecto de cociente \(129\). Suprimirlo cambiaría la aritmética que se transporta, aunque la tabla residual permaneciera idéntica.

\subsubsection{Reducción ternaria y radical multiplicativo}

El anillo \(R=\ZZ/9\ZZ\) tiene grupo de unidades \(\{1,2,4,5,7,8\}\) e ideal máximo
\[
 \mathfrak m=(3)=\{0,3,6\},\qquad\mathfrak m^2=0.
\]
Toda fila aditiva es una permutación de los nueve representantes. Una fila multiplicativa lo es exactamente cuando su índice es una unidad. Las filas tres y seis contienen tres valores repetidos; la fila nueve contiene únicamente el representante de la clase nula. Por ello el plegamiento multiplicativo ya distingue, antes de cualquier geometría continua, las unidades del sector nilpotente.

La aplicación \(x\mapsto3x\) tiene núcleo e imagen iguales a \(\mathfrak m\). Induce un isomorfismo de espacios vectoriales sobre \(\FF_3\),
\[
 R/\mathfrak m\longrightarrow\mathfrak m,\qquad [x]\longmapsto3x,
\]
pero no un isomorfismo de anillos. En particular, la secuencia visible \(3,6,9\) admite dos lecturas diferentes: su reducción directa módulo tres es \(0,0,0\), mientras que la extracción del factor tres seguida de la reducción de su coordenada interna da \(1,2,0\). Esta segunda operación conserva una coordenada del ideal; no es la misma proyección que reducir directamente el número original.

\subsubsection{Realización latina de la evaluación aditiva}

La diferencia entre las dos evaluaciones puede expresarse también mediante
operadores. Sea \((e_r)_{r\in\ZZ/9\ZZ}\) la base ortonormal de \(\CC^9\).
La asignación \(v_{ab}=e_{a+b}\) produce una base ortonormal en cada fila
y columna. Los proyectores \(P_{ab}=v_{ab}v_{ab}^{*}\) satisfacen
\[
 \sum_bP_{ab}=I_9=\sum_aP_{ab},\qquad P_{ab}^2=P_{ab}.
\]
En efecto, cada traslación \(b\mapsto a+b\) permuta los nueve residuos.
Se obtiene una realización por proyectores de rango uno de un cuadrado
latino, en la terminología de los cuadrados latinos cuánticos
\cite{x_reciprocidad:mustovicary,x_reciprocidad:delascuevas2023}. Esta construcción parte de la hoja
aditiva ya definida y explicita el mapa hacia su representación operatoria.

La evaluación multiplicativa conserva otro contenido: para \(a=3\), los
vectores \(e_{3b}\) repiten los residuos \(0,3,6\), y
\[
 \sum_{b\in\ZZ/9\ZZ}|e_{3b}\rangle\langle e_{3b}|
 =3\bigl(|e_0\rangle\langle e_0|+|e_3\rangle\langle e_3|
                    +|e_6\rangle\langle e_6|\bigr).
\]
La multiplicidad queda registrada en el operador y caracteriza el radical
multiplicativo. Por ello el soporte toroidal, la propiedad latina y la
resolución operatoria de la identidad son propiedades que requieren
definiciones diferentes. Los cuadrados mágicos cuánticos estudiados por
De las Cuevas, Drescher y Netzer exigen entradas positivas y sumas de filas
y columnas iguales a la identidad; su teoría distingue además las
dilataciones con entradas conmutativas\cite{x_reciprocidad:delascuevas2020}. Esta referencia
proporciona un lenguaje preciso para comparar realizaciones de APP.

\subsection{TRIT: orientación y regímenes cuadráticos}

\begin{definition}[Firma ternaria y levantamiento local]
La firma TRIT utiliza la identificación balanceada
\[
 \FF_3\longrightarrow\TRIT,\qquad 0\mapsto0,\quad1\mapsto+1,\quad2\mapsto-1.
\]
Para \(n\in\ZZ\), su levantamiento conserva \(n=r+3c\), donde \(r\in\TRIT\) y \(c\in\ZZ\). En una transición de amplitud local acotada, los tres estados \((0,-1),(0,0),(0,+1)\) pueden presentar el mismo residuo nulo y distintos acarreos.
\end{definition}

La firma no reemplaza a la trayectoria. Un cero residual no significa ausencia de información de orientación, fase o cociente. Tampoco introduce por sí solo una magnitud física. Su función es distinguir algebraicamente las clases de transporte que se realizarán a continuación.

\subsubsection{Tres regímenes cuadráticos y una exponencial común}

Una vez obtenida la firma ternaria, su realización cuadrática se define por
\begin{equation}
 \mathbb A_\tau=\RR[X]/(X^2+\tau),\qquad \tau\in\TRIT.
 \label{x_reciprocidad:eq:algebras-trit}
\end{equation}
Si \(J_\tau\) es la clase de \(X\), entonces \(J_\tau^2=-\tau I\). La estructura distingue el tipo elíptico \(\tau=+1\), el parabólico \(\tau=0\) y el hiperbólico \(\tau=-1\). La realización real no genera la firma: expresa los tipos que la reducción orientada ya ha producido.

\begin{proposition}[Exponenciación de los tres tipos]
En cada álgebra \eqref{x_reciprocidad:eq:algebras-trit},
\begin{equation}
 \exp(tJ_\tau)=C_\tau(t)I+S_\tau(t)J_\tau,
\end{equation}
donde
\[
 C_\tau(t)=\sum_{n\ge0}\frac{(-\tau)^nt^{2n}}{(2n)!},\qquad
 S_\tau(t)=\sum_{n\ge0}\frac{(-\tau)^nt^{2n+1}}{(2n+1)!}.
\]
Las tres realizaciones son
\[
 \begin{array}{c|c|c}
 \tau&J_\tau^2&\exp(tJ_\tau)\\\hline
 +1&-I&\cos t\,I+\sin t\,J_\tau\\
 0&0&I+tJ_\tau\\
 -1&I&\cosh t\,I+\sinh t\,J_\tau.
 \end{array}
\]
\end{proposition}
\begin{proof}
Separar los términos pares e impares de la serie exponencial y sustituir \(J_\tau^{2n}=(-\tau)^nI\) y \(J_\tau^{2n+1}=(-\tau)^nJ_\tau\) da las fórmulas. Para \(\tau=0\), todos los términos de grado al menos dos se anulan. En los otros dos casos se recuperan las series trigonométricas o hiperbólicas.
\end{proof}

La exponencial es común a los tres regímenes: no se asigna exclusivamente a la rama parabólica. Los invariantes
\[
 C_\tau(t)^2+\tau S_\tau(t)^2=1
\]
se deducen de \(\exp(tJ_\tau)\exp(-tJ_\tau)=I\). El mismo principio algebraico conserva la unidad a través de realizaciones que difieren por su tipo cuadrático.

En el régimen hiperbólico, \(P_\pm=(I\pm J_{-1})/2\) son proyectores complementarios y
\[
 \exp(tJ_{-1})=e^tP_++e^{-t}P_-.
\]
El crecimiento y la contracción aparecen conjuntamente en dos direcciones propias. En el elíptico, el parámetro describe una rotación; en el parabólico, una traslación nilpotente. La interpretación temporal adoptada por HMT asocia \(+1\) al retorno con memoria, \(0\) al umbral y \(-1\) a la apertura bidireccional. Esta interpretación exige un orden de parametrización; no añade un cuarto tipo algebraico.


% BEGIN OWNER /Users/ruben/Documents/New project/output/EXTREMA_MEDIA_RAZON_RECIPROCIDAD_ES_EN_20260918/ARTICULO/ES/source/sections/trit_desarrollo.tex
% Ampliación acumulativa del TRIT; se inserta antes de la figura de regímenes.
% No sustituye ni elimina el desarrollo precedente.
\subsubsection{División balanceada y composición de los levantamientos}
\label{x_reciprocidad:trit-ampl-div-balanceada}

La reducción ternaria permite expresar una evaluación entera de APP en una
coordenada visible y una coordenada de prolongación. Su función no consiste
únicamente en abreviar la escritura de un entero: proporciona una regla de
composición que determina cómo se modifica el cociente cuando se suman o se
multiplican evaluaciones. La elección balanceada hace explícita, además, la
compatibilidad de esa regla con la inversión de signo.

\begin{definition}[Coordenadas ternarias balanceadas]
\label{x_reciprocidad:trit-ampl-def-coordenadas}
Para \(n\in\mathbb Z\), sean
\[
 q_3(n)=\left\lfloor\frac{n+1}{3}\right\rfloor,\qquad
 r_3(n)=n-3q_3(n).
\]
La aplicación
\[
 \mathcal B:\mathbb Z\longrightarrow\TRIT\times\mathbb Z,\qquad
 n\longmapsto(r_3(n),q_3(n))
\]
se denomina descomposición ternaria balanceada. Su inversa es
\(\mathcal L(r,c)=r+3c\).
\end{definition}

En efecto, \(r_3(n)\in\{-1,0,+1\}\). Si dos pares representan el mismo entero,
\(r+3c=s+3d\), entonces \(r-s\) es múltiplo de \(3\) y pertenece al intervalo
\([-2,2]\); necesariamente \(r=s\) y \(c=d\). La representación es, por tanto,
única. La división puede repetirse sobre \(q_3(n)\), de manera que cada nivel
conserva el dato necesario para reconstruir el anterior. El cociente no es una
perturbación del residuo: constituye su coordenada complementaria.

La firma obtenida de una evaluación \(F\), con \(F=S\) o \(F=P\) en APP,
es \(r_3\circ F\). Cuando la reducción actúa sobre el cursor de fase, la misma
sección balanceada produce las clases \(\{1,4,7\}\), \(\{2,5,8\}\) y
\(\{3,6,9\}\), con valores \(+1,-1,0\). El selector del TPK que se define
a continuación utiliza esas clases para determinar el modo operativo.
La lectura de una evaluación y la selección de un modo tienen así la misma
reducción ternaria, pero dominios y funciones diferentes.

\Needspace{16\baselineskip}
\begin{proposition}[Aritmética de los pares balanceados]
\label{x_reciprocidad:trit-ampl-prop-aritmetica}
Sean \(n=r+3c\) y \(m=s+3d\), con \(r,s\in\TRIT\). Defínase
\[
 \chi(r,s)=\frac{r+s-r_3(r+s)}3.
\]
Las operaciones enteras se transportan a los pares mediante
\begin{align}
 (r,c)\boxplus(s,d)
 &=\bigl(r_3(r+s),c+d+\chi(r,s)\bigr),
 \label{x_reciprocidad:trit-ampl-eq-suma}\\
 (r,c)\boxtimes(s,d)
 &=\bigl(rs,rd+sc+3cd\bigr).
 \label{x_reciprocidad:trit-ampl-eq-producto}
\end{align}
La aplicación \(\mathcal L\) conserva ambas operaciones.
\end{proposition}
\begin{proof}
La igualdad \(r+s=r_3(r+s)+3\chi(r,s)\) da
\[
 n+m=r_3(r+s)+3\bigl(c+d+\chi(r,s)\bigr).
\]
Por otra parte,
\[
 nm=rs+3(rd+sc+3cd).
\]
Como el producto de dos representantes balanceados pertenece de nuevo a
\(\{-1,0,+1\}\), no necesita una reducción residual adicional. La unicidad
de la descomposición demuestra las dos fórmulas.
\end{proof}

Por ejemplo, \(2=(-1)+3\) y \(4=1+3\). Su suma y su producto se reconstruyen
como
\[
 (-1,1)\boxplus(1,1)=(0,2),\qquad
 (-1,1)\boxtimes(1,1)=(-1,3).
\]
Las evaluaciones resultantes son \(6\) y \(8\). La suma de cocientes basta
en este ejemplo aditivo; en el producto intervienen los términos cruzados
\(rd\), \(sc\) y \(3cd\). La discrepancia entre ambas operaciones reside
también en la coordenada de prolongación, no solamente en el residuo.

El término \(\chi\) es un cociclo de la sección balanceada. Si
\(r\oplus s=r_3(r+s)\), satisface
\[
 \chi(r,s)+\chi(r\oplus s,t)
 =\chi(s,t)+\chi(r,s\oplus t).
\]
Ambos miembros representan el cociente de la misma suma de tres enteros.
Esta identidad garantiza que reagrupar las operaciones no cambia la
evaluación reconstruida. Al mismo tiempo, la proyección sobre el primer
componente suprime esa contabilidad: \(1+1=-1+3\), mientras que la lectura
residual conserva únicamente \(-1\). La información eliminada está
determinada, en este caso, por \(\chi(1,1)=1\).

La relación con el módulo \(9\) conserva una segunda escala discreta.
La aplicación
\[
 \beta:\TRIT\times\mathbb F_3\longrightarrow\mathbb Z/9\mathbb Z,
 \qquad (r,\bar c)\longmapsto r+3c\pmod9
\]
está bien definida y es biyectiva. Cambiar el representante de \(\bar c\)
añade un múltiplo de \(9\); recíprocamente, una igualdad de imágenes fija
primero \(r\) por reducción módulo \(3\) y después \(\bar c\). La suma
transportada por esta biyección contiene el término
\(\chi(r,s)\) de \eqref{x_reciprocidad:trit-ampl-eq-suma}, reducido ahora módulo \(3\).
Así, los dos niveles ternarios no se componen mediante una suma
independiente de coordenadas: la normalización del primer nivel modifica
el segundo.

En particular, las clases \(3\), \(6\) y \(0\) de
\(\mathbb Z/9\mathbb Z\) tienen primer componente \(0\), pero segundo
componente \(1\), \(2\) y \(0\), respectivamente. La distinción ya observada
entre reducción directa y extracción de la coordenada del ideal aparece
como una propiedad de \(\beta\). Conservar esa coordenada impide que el
sector residual nulo se convierta artificialmente en una sola posición.
El módulo \(1000\), utilizado después para organizar bloques decimales,
cumple otra función de representación posicional; su cociente se conserva
mediante la división correspondiente. Un cambio de base conserva el entero
cuando retiene residuo y cociente, mientras que la coincidencia de residuos
en bases distintas no identifica por sí sola sus estados de transporte.

\Needspace{11\baselineskip}
\subsubsection{Inversión, clase residual y coordenadas no visibles}
\label{x_reciprocidad:trit-ampl-inversion}

La inversión aditiva se levanta sin ambigüedad:
\[
 \mathcal B(-n)=(-r_3(n),-q_3(n)).
\]
En particular, \(3\), \(0\) y \(-3\) tienen coordenadas
\((0,1)\), \((0,0)\) y \((0,-1)\). Comparten firma nula, pero representan
evaluaciones distintas. En una transición restringida a un cociente de
amplitud máxima \(1\), son las tres posibilidades de la fibra visible sobre
cero; en un levantamiento entero general, esa fibra contiene todos los pares
\((0,c)\), con \(c\in\mathbb Z\).

La coordenada de cociente recupera el entero evaluado, pero no recupera por
sí sola la trayectoria que lo produjo. Dos caminos de APP pueden compartir
evaluación, residuo y cociente, y diferir en orientación, sucesión de modos o
enrollamiento. Por ello la descomposición \(\mathcal B\) se integra en el
estado de transporte, en vez de sustituirlo. La memoria aritmética de una
división y la memoria de una trayectoria son coordenadas relacionadas, no
sinónimos.

Para una sucesión de firmas de longitud \(n\), la inversión orientada es
\[
 C_{\rm rev}(\tau_1,\ldots,\tau_n)
   =(-\tau_n,\ldots,-\tau_1).
\]
Aplicarla dos veces restituye la sucesión. La operación invierte el orden de
las posiciones y el signo de cada firma; conserva su longitud y las
posiciones nulas con el orden invertido. En la parametrización adoptada,
intercambia el retorno con memoria, asociado a \(+1\), y la apertura
bidireccional, asociada a \(-1\), manteniendo el umbral \(0\). El transporte
de los restantes datos de una trayectoria se especifica mediante la
involución correspondiente del TPK.

Esta inversión de firmas debe distinguirse de la conjugación dentro de un
álgebra cuadrática. La primera puede cambiar el tipo \(\tau\); la segunda,
que se construye a continuación, actúa dentro del tipo ya fijado. La
distinción impide atribuir a una inversión de signo una equivalencia
algebraica entre regímenes diferentes.

\subsubsection{Producto cuadrático, conjugación y norma algebraica}
\label{x_reciprocidad:trit-ampl-norma}

Un elemento de \(\mathbb A_\tau\) tiene expresión única \(aI+bJ_\tau\).
La relación \(J_\tau^2=-\tau I\) determina su producto:
\begin{equation}
 (aI+bJ_\tau)(cI+dJ_\tau)
   =(ac-\tau bd)I+(ad+bc)J_\tau.
 \label{x_reciprocidad:trit-ampl-producto-cuadratico}
\end{equation}
Se trata de una realización real del tipo ya seleccionado. Las coordenadas
\(a,b\) no son residuos de APP, y el producto
\eqref{x_reciprocidad:trit-ampl-producto-cuadratico} no sustituye a
\eqref{x_reciprocidad:trit-ampl-eq-producto}; expresa otra operación, en el codominio
cuadrático declarado.

\begin{definition}[Conjugación y norma algebraica]
\label{x_reciprocidad:trit-ampl-def-norma}
La conjugación del tipo \(\tau\) y su norma algebraica son
\[
 \overline{aI+bJ_\tau}=aI-bJ_\tau,\qquad
 N_\tau(aI+bJ_\tau)=a^2+\tau b^2.
\]
\end{definition}

\begin{proposition}[Multiplicatividad e invertibilidad]
\label{x_reciprocidad:trit-ampl-prop-norma}
Para \(z,w\in\mathbb A_\tau\),
\[
 z\bar z=N_\tau(z)I,\qquad
 N_\tau(zw)=N_\tau(z)N_\tau(w).
\]
Además, \(z\) es invertible si y sólo si \(N_\tau(z)\ne0\), y en ese caso
\[
 z^{-1}=\frac{\bar z}{N_\tau(z)}.
\]
\end{proposition}
\begin{proof}
El producto de \(aI+bJ_\tau\) y \(aI-bJ_\tau\) es
\((a^2+\tau b^2)I\). La conjugación es multiplicativa y el álgebra es
conmutativa, por lo que
\((zw)\overline{zw}=(z\bar z)(w\bar w)\). Si la norma no se anula, la
fórmula indicada proporciona la inversa. Recíprocamente, si \(zw=I\), la
multiplicatividad implica \(N_\tau(z)N_\tau(w)=1\).
\end{proof}

La forma cuadrática \(N_\tau\) es positiva definida únicamente en el tipo elíptico.
Para \(\tau=0\), vale \(a^2\), y el elemento \(J_0\) es no nulo aunque su
cuadrado y su norma sean cero. Para \(\tau=-1\), vale \(a^2-b^2\); los
elementos no nulos \(I+J_{-1}\) e \(I-J_{-1}\) tienen norma cero y su
producto se anula. La norma algebraica es aquí una forma cuadrática
multiplicativa: no presupone una norma positiva de espacio vectorial.

Los tres tipos pueden realizarse fielmente mediante
\[
 aI+bJ_\tau\longmapsto
 \begin{pmatrix}a&-\tau b\\b&a\end{pmatrix}.
\]
El determinante de esta matriz es \(N_\tau(aI+bJ_\tau)\). Quedan
representadas así, con una misma expresión matricial, la estructura
compleja, la estructura de números duales y la estructura real escindida.
Cada una conserva su ley y sus divisores de cero. La forma de firma
\((1,1)\) del último caso tiene interpretación lorentziana en dimensión
\(1+1\); la identificación de una métrica física requiere además un espacio,
un campo y una regla de transporte específicos.

\subsubsection{Conservación multiplicativa y composición de evoluciones}
\label{x_reciprocidad:trit-ampl-evoluciones}

La exponencial ya obtenida pertenece al conjunto de elementos de norma uno:
\[
 N_\tau\bigl(\exp(tJ_\tau)\bigr)=1.
\]
Esta conservación se mantiene al componer parámetros. En efecto, todos los
factores son funciones del mismo generador, y
\[
 \exp(tJ_\tau)\exp(sJ_\tau)=\exp((t+s)J_\tau).
\]
Al comparar las dos componentes se deducen las leyes de adición
\begin{align}
 C_\tau(t+s)&=C_\tau(t)C_\tau(s)-\tau S_\tau(t)S_\tau(s),\\
 S_\tau(t+s)&=S_\tau(t)C_\tau(s)+C_\tau(t)S_\tau(s).
\end{align}
La conjugación cambia \(t\) por \(-t\) y coincide con la inversión dentro
de este subgrupo. No cambia \(\tau\): invierte una evolución del régimen
elegido.

En el tipo elíptico, la forma conservada es \(a^2+b^2\), y la trayectoria
exponencial recorre el círculo de norma uno. En el hiperbólico, las
coordenadas propias son \(u=a+b\), \(v=a-b\), con \(uv=N_{-1}(z)\).
La exponencial da
\[
 (u,v)=(e^t,e^{-t}),\qquad uv=1.
\]
El crecimiento de una coordenada se acompaña de la contracción de la otra.
La conservación expresa una relación entre ambas coordenadas; no afirma que
cada una permanezca constante.

En el tipo parabólico,
\[
 (I+tJ_0)(I+sJ_0)=I+(t+s)J_0,\qquad
 (I+tJ_0)^{-1}=I-tJ_0.
\]
La evolución no se detiene porque \(J_0^2=0\). La nilpotencia elimina los
términos de orden al menos dos, pero conserva el término lineal que
transporta el parámetro. Son distintos el símbolo visible \(0\), el
elemento nulo del álgebra y el generador no nulo \(J_0\). Esta distinción
es necesaria para describir el umbral sin identificarlo con ausencia de
estado o de dinámica.

\begin{proposition}[Composición de coordenadas afines]
\label{x_reciprocidad:trit-ampl-prop-pendientes}
En la carta donde la componente escalar no se anula, represéntese una
dirección mediante \(I+uJ_\tau\). Si \(1-\tau uv\ne0\), la dirección de su
producto con \(I+vJ_\tau\) tiene coordenada
\[
 u\star_\tau v=\frac{u+v}{1-\tau uv}.
\]
\end{proposition}
\begin{proof}
Se tiene
\[
 (I+uJ_\tau)(I+vJ_\tau)
   =(1-\tau uv)I+(u+v)J_\tau.
\]
Dividir por la componente escalar da la expresión anunciada.
\end{proof}

La carta afín conserva la razón entre componentes, no el factor escalar
dividido. Cuando \(1-\tau uv=0\), debe examinarse el producto antes de
cambiar de carta: puede tener dirección definida con componente escalar
nula, o ser el elemento cero en presencia de divisores de cero. La ley de
composición sigue perteneciendo al álgebra completa; la singularidad de
una coordenada no equivale automáticamente a una singularidad del objeto.
Esta fórmula reaparecerá en la lectura regional elíptica. Su uso requiere
mantener el denominador y, cuando corresponda, el factor de normalización.

\subsubsection{Transporte del tipo y realizaciones posteriores}
\label{x_reciprocidad:trit-ampl-puentes}

Las álgebras \(\mathbb A_{+1}\), \(\mathbb A_0\) y
\(\mathbb A_{-1}\) no son isomorfas como álgebras reales unitarias.
La primera es un cuerpo; la segunda contiene un nilpotente no nulo; la
tercera posee idempotentes no triviales y es isomorfa a
\(\mathbb R\times\mathbb R\). Por ello, el cambio de firma que organiza una
ruta no puede interpretarse como un isomorfismo indiferenciado entre las
tres realizaciones. El TPK conserva el índice de régimen, los dominios y el
operador efectivo que actúa en cada transición.

Dentro del capítulo electrónico, el par de proyecciones de las hojas de APP
produce un conmutador normalizado cuyo cuadrado es \(-I\). El mismo bloque
contiene una involución de cuadrado \(I\) y dos direcciones nilpotentes. La
realización conjunta se demostrará sobre ese par concreto de proyecciones;
la clasificación precedente explica por qué las tres relaciones cuadráticas
pueden aparecer en una misma álgebra matricial sin identificarse entre sí.

En la prolongación excepcional, la matriz de Paley reducida módulo \(3\)
satisface \(A_W^2=-I_6\). La relación dota al espacio ternario de una
estructura de dimensión \(3\) sobre \(\mathbb F_9\). Su vínculo con una raíz
operatorial real se expresa mediante la presentación integral
\[
 \mathbb Z[u]/(u^2+1).
\]
Los cambios de base a \(\mathbb F_3\) y a \(\mathbb R\) producen
representaciones de esa misma relación. Cada una conserva su característica,
su dimensión y su información de incidencia. El desarrollo de los
respectivos operadores establece el enlace, no la igualdad de sus
cardinales.

La conservación de norma uno describe, por consiguiente, las realizaciones
cuadráticas de las evoluciones. La conservación de cocientes, orientación y
trayectoria pertenece al estado aritmético transportado. Ambas propiedades
intervienen en la construcción, con funciones diferentes: una determina las
identidades de las representaciones; la otra permite componer, reconstruir y
prolongar las operaciones que las originan.

% END OWNER



% BEGIN OWNER /Users/ruben/Documents/New project/output/EXTREMA_MEDIA_RAZON_RECIPROCIDAD_ES_EN_20260918/ARTICULO/ES/source/figures/regimenes_trit.tex
\begin{figure}[htbp]
\centering
\begin{tikzpicture}[scale=.88,>=Stealth,font=\small]
\foreach \c/\titulo in {0/{Elíptico},4.6/{Parabólico},9.2/{Hiperbólico}}{
\begin{scope}[xshift=\c cm]
\draw[->,gray] (-1.65,0)--(1.7,0) node[right] {\(u\)};
\draw[->,gray] (0,-1.35)--(0,1.5) node[above] {\(v\)};
\node at (0,2.05) {\titulo};
\end{scope}}
\draw[thick] (0,0) circle[radius=1];
\draw[->,thick] (1,0) arc[start angle=0,end angle=60,radius=1];
\node at (0,-1.85) {\(u^2+v^2=1\)};
\node at (0,-2.35) {\(J^2=-I\)};
\draw[thick] (3.6,-1.2)--(3.6,1.2);
\draw[->,thick] (5.6,-1.2)--(5.6,1.2);
\node at (4.6,-1.85) {\(u^2=1\)};
\node at (4.6,-2.35) {\(J^2=0\)};
\draw[thick,domain=-.9:.9,samples=50] plot ({9.2+cosh(\x)},{sinh(\x)});
\draw[thick,domain=-.9:.9,samples=50] plot ({9.2-cosh(\x)},{sinh(\x)});
\node at (9.2,-1.85) {\(u^2-v^2=1\)};
\node at (9.2,-2.35) {\(J^2=I\)};
\end{tikzpicture}
\caption{Realizaciones cuadráticas de TRIT. Para
\(J_\tau=\left(\begin{smallmatrix}0&-\tau\\1&0\end{smallmatrix}\right)\),
la exponencial conserva \(u^2+\tau v^2\). La degeneración parabólica
conserva \(u\) y transporta \(v\) linealmente. Los tres diagramas representan
órbitas de formas cuadráticas; su tipo geométrico corresponde a la relación
algebraica del generador.}
\label{x_reciprocidad:fig:trit-regimenes}
\end{figure}

% END OWNER


\newpage
\subsection{TPK: núcleo topológico de fase}

La dinámica utiliza el selector de fase
\[
 M_{\rm ph}(r)=
 \begin{cases}
 +1,&r\in\{1,4,7\},\\
 -1,&r\in\{2,5,8\},\\
 0,&r\in\{3,6,9\}.
 \end{cases}
\]
Su valor decide qué evaluación opera. El transporte cardinal desplaza el cursor correspondiente; el registro conserva el valor anterior, el nuevo cociente, la orientación y la transición. Por tanto, selección y transporte son funciones distintas: un signo ternario no es un desplazamiento del toro.

Denotaremos por \(\widehat X\) el espacio de estados aritméticos con memoria de trayectoria. Un elemento contiene las posiciones y direcciones de los dos cursores, el modo activo, la firma TRIT, la fase, los residuos y cocientes de ambas evaluaciones, los acarreos, el prefijo construido y los datos de frontera. Cuando se introduce un cilindro de refinamiento, éste forma parte del estado. La notación evita sustituir el objeto por una versión reducida que conserve únicamente el símbolo visible.

\begin{definition}[Actualización del TPK]
Sean \(\mathsf{Sel}_t\) la selección de modo mediante \(M_{\rm ph}\), \(\mathsf{Tra}_t\) el transporte cardinal levantado y \(\mathsf{Mem}_t\) la actualización de los registros. La transición es
\begin{equation}
 \mathcal U_t=\mathsf{Mem}_t\circ\mathsf{Tra}_t\circ\mathsf{Sel}_t.
 \label{x_reciprocidad:eq:actualizacion}
\end{equation}
Cada aplicación conserva los campos que no modifica y actualiza los restantes mediante la división euclídea y la concatenación de la trayectoria.
\end{definition}

La notación \eqref{x_reciprocidad:eq:actualizacion} expresa el orden de composición, no reemplaza las reglas efectivas. La siguiente sección especifica completamente su realización finita de dos cursores. Después, la prolongación no se obtendrá borrando su memoria y repitiendo la primera emisión, sino transportando los prefijos y sus fronteras compatibles. De este modo, APP proporciona las evaluaciones, TRIT distingue sus regímenes y TPK compone su evolución.


% BEGIN OWNER /Users/ruben/Documents/New project/output/EXTREMA_MEDIA_RAZON_RECIPROCIDAD_ES_EN_20260918/ARTICULO/ES/source/sections/tpk_desarrollo_integrado.tex
% Recuperación expositiva desde los propietarios del núcleo TPK.
% La procedencia y el control de conservación se documentan en technical/TPK_INTEGRACION.md.
% Este archivo amplía la sección TPK; no sustituye extension.tex ni generacion.tex.

\subsubsection{Estado dinámico, observables y dependencia de las operaciones}
\label{x_reciprocidad:tpk-int:estado}

El núcleo topológico de fase organiza la evolución de las dos hojas de APP
con la orientación determinada por TRIT. Su definición comprende el estado
sobre el que actúa, las reglas de transición y las relaciones que hacen
compatibles las prolongaciones. La ecuación \eqref{x_reciprocidad:eq:actualizacion} adquiere
contenido operativo al especificar qué coordenadas recibe y modifica cada
uno de sus factores. La posición pertenece al soporte toroidal; la fase
pertenece al calendario; la orientación determina el transporte; los
cocientes registran la evaluación aritmética anterior a su reducción.
Estas coordenadas se relacionan mediante la transición y conservan funciones
matemáticas diferentes.

La proyección observable de un estado es
\begin{equation}
 q=(x_+,d_+,x_\times,d_\times,\theta)
 \in\mathcal Q_{\rm obs}
 :=(X_9\times\mathcal D)^2\times\Theta_{108},
 \qquad \mathcal D=\{N,E,S,O\},\quad
 \Theta_{108}=\ZZ/108\ZZ.
 \label{x_reciprocidad:tpk-int:qobs}
\end{equation}
Los subíndices distinguen los cursores aditivo y multiplicativo. En las
fibras de estados completos sobre \(q\) se conservan la trayectoria ordenada, las evaluaciones
enteras de ambas hojas, sus residuos y cocientes, y los registros de
orientación y frontera. En una prolongación de precisión se añaden el
prefijo y el cilindro compatible. Por tanto, dos estados con la misma
proyección \(q\) pueden conservar registros distintos.

La dependencia de las operaciones tiene tres niveles. En el nivel local,
el selector activa una hoja, el transporte modifica su posición y la
actualización calcula las nuevas coordenadas aritméticas. En el nivel de
trayectorias, estas transiciones se componen y conservan los registros de
los pasos anteriores. En el nivel de refinamiento, las relaciones de
extensión determinan qué continuación conserva el prefijo, la firma y la
frontera del antecedente. El calendario finito describe la fase de esas
operaciones; la holonomía describe su composición sobre la historia.

\subsubsection{Selección trítica y transporte cardinal levantado}
\label{x_reciprocidad:tpk-int:seleccion}

Sea \(\bar\theta\in\{0,\ldots,107\}\) el representante del calendario.
La fase operativa y el sector dodecafásico son
\begin{equation}
 r(\theta)=1+(\bar\theta\bmod9),\qquad
 \beta(\theta)=\left[\left\lfloor\frac{\bar\theta}{9}\right\rfloor\right]_{12}.
 \label{x_reciprocidad:tpk-int:fase}
\end{equation}
El selector ya definido puede escribirse \(\varepsilon_9=M_{\rm ph}\).
Su vector de valores, en la carta ordenada de \(D_9\), es
\[
 m_{\rm ph}=(1,-1,0,1,-1,0,1,-1,0).
\]
La función \(M_{\rm ph}:D_9\to\TRIT\) y este vector son dos
presentaciones del mismo selector, con sus dominios respectivos.

Definamos los indicadores de activación
\[
 a_+(\theta)=\mathbf1_{\{1,4,7\}}(r(\theta)),\qquad
 a_\times(\theta)=\mathbf1_{\{2,5,8\}}(r(\theta)).
\]
El transporte de los cursores es entonces
\begin{equation}
 x_+'=x_++a_+(\theta)\nu_{d_+},\qquad
 x_\times'=x_\times+a_\times(\theta)\nu_{d_\times}
 \quad\text{en }X_9.
 \label{x_reciprocidad:tpk-int:cursores}
\end{equation}
En las fases \(3,6,9\) permanecen ambas posiciones; el evento neutral
forma parte de la cronología y de su registro. Las direcciones cardinales
determinan los vectores \(\nu_d\), mientras la firma de fase determina
cuál de esos vectores se aplica.

El levantamiento entero conserva el cruce de frontera. Si
\(s:X_9\to D_9^2\) elige los representantes positivos y
\(x'=x+a\nu_d\), con \(a\in\{0,1\}\), se define
\begin{equation}
 b_d(x,a)=\frac{s(x)+a\nu_d-s(x')}{9}\in\ZZ^2.
 \label{x_reciprocidad:tpk-int:borde}
\end{equation}
La divisibilidad por \(9\) se sigue de la igualdad en \(X_9\).
Si \(z\in\ZZ^2\) registra los cruces anteriores, la actualización es
\(z'=z+b_d(x,a)\). De este modo,
\[
 s(x')+9z'=s(x)+9z+a\nu_d.
\]
El dato entero de transporte permanece reconstruible después de reducir
la posición sobre el toro.

\begin{proposition}[Conservación del desplazamiento levantado]
\label{x_reciprocidad:tpk-int:prop-desplazamiento}
Para una trayectoria cardinal de \(n\) pasos, cuyos indicadores de
activación son \(a_t\), se cumple
\[
 s(x_n)+9z_n=s(x_0)+9z_0+
 \sum_{t=0}^{n-1}a_t\nu_{d_t}.
\]
La misma igualdad se conserva al concatenar dos trayectorias compatibles.
\end{proposition}
\begin{proof}
La identidad de un paso es \eqref{x_reciprocidad:tpk-int:borde}. Al sumarla para los
\(n\) pasos, las coordenadas intermedias se cancelan. Dividir el intervalo
de suma en dos tramos produce exactamente la ley de concatenación. En una
trayectoria cerrada, la diferencia \(z_n-z_0\) es el grado de
enrollamiento definido previamente.
\end{proof}

El mismo principio rige las evaluaciones de APP. Si una hoja toma el
valor entero \(N_t\), su registro conserva simultáneamente
\(N_t=R_t+9q_t\). La diferencia entre dos pasos verifica
\[
 N_{t+1}-N_t=(R_{t+1}-R_t)+9(q_{t+1}-q_t).
\]
La memoria de cruce espacial \(z_t\) y el cociente aritmético \(q_t\)
tienen definiciones diferentes; su conservación conjunta mantiene la
procedencia geométrica y aritmética de cada transición.

\subsubsection{Cronología de dos cursores y retornos compuestos}
\label{x_reciprocidad:tpk-int:cronologia}

La involución cardinal \(\jmath\) intercambia \(N\) con \(S\) y
\(E\) con \(O\). Tras aplicar \eqref{x_reciprocidad:tpk-int:cursores}, ambas
direcciones se invierten cuando finaliza un bloque de \(27\) pasos.
Escribiendo
\[
 h(\theta)=\mathbf1_{\{0,27,54,81\}}
             ((\bar\theta+1)\bmod108),
\]
la regla completa sobre la proyección observable es
\begin{equation}
 F_{\rm obs}(q)=
 \bigl(x_+',\jmath^{h(\theta)}d_+,
       x_\times',\jmath^{h(\theta)}d_\times,
       \theta+1\bigr).
 \label{x_reciprocidad:tpk-int:fobs}
\end{equation}
La lectura que forma la emisión utiliza las posiciones anteriores al
transporte, según la recurrencia explícita de la
sección~\ref{x_reciprocidad:sec:extension}. La actualización observable y la emisión
se componen, por tanto, con un orden fijado.

\begin{proposition}[Reversibilidad y periodo del calendario observable]
\label{x_reciprocidad:tpk-int:periodo}
La aplicación \(F_{\rm obs}\) es una permutación de
\(\mathcal Q_{\rm obs}\) y satisface \(F_{\rm obs}^{108}=I\).
En un bloque de \(54\) pasos se restauran posiciones y direcciones;
la coordenada de \(\Theta_{108}\) avanza \(54\).
\end{proposition}
\begin{proof}
Dado el estado final, se recupera primero la fase anterior restando uno
en \(\Theta_{108}\). Esa fase determina si se invirtieron las
direcciones y qué cursor se desplazó. Aplicar la involución correspondiente
y restar el vector cardinal recupera el estado inicial de manera única.

Cada intervalo de \(27\) pasos contiene nueve activaciones de cada
cursor. El bloque siguiente conserva las activaciones e invierte las
direcciones, de modo que los desplazamientos enteros se cancelan en
\(54\) pasos. La regla tiene periodo \(54\) en sus incrementos de
posición y orientación; la misma cancelación vale con cualquier origen
del calendario. Tras \(108\) pasos se restaura también \(\theta\).
\end{proof}

Se reserva
\[
 \mathcal K_{27}=F_{\rm obs}^{27}:
 \mathcal Q_{\rm obs}\longrightarrow\mathcal Q_{\rm obs}
\]
para el iterado de \(27\) transiciones. Su cuarto iterado es la
identidad sobre el estado observable definido en \eqref{x_reciprocidad:tpk-int:qobs}.
La proyección de la historia al calendario permite efectuar esta
comprobación finita. El registro de la trayectoria, sus evaluaciones y
los nuevos prefijos se conservan en el estado dinámico sobre el que
actúa \(\mathcal U_t\); la igualdad observable no los reinicializa.

\subsubsection{Transporte afín de la memoria y composición de estados}
\label{x_reciprocidad:tpk-int:memoria}

La proyección observable admite fibras de datos aritméticos y geométricos.
Sobre una configuración \(q\), escribimos un estado en la forma
\[
 x=(q,o,h,c,\sigma,C,b,\lambda).
\]
Aquí \(o\) conserva la orientación y la hoja; \(h\), los residuos
y cocientes; \(c\), la memoria aditiva de transporte; \(\sigma\),
la firma de frontera; \(C\), el cilindro de precisión; \(b\), su
etiqueta de frontera; y \(\lambda\), la trayectoria registrada. La
firma está determinada por una aplicación
\[
 \operatorname{Sig}_q:
 \mathsf H_q\times\mathsf C_q\times\mathsf{Hist}_q
 \longrightarrow\mathsf S_q,
 \qquad \sigma=\operatorname{Sig}_q(h,c,\lambda).
\]
Los símbolos \(\mathsf H_q,\mathsf C_q,\mathsf{Hist}_q\) y
\(\mathsf S_q\) denotan, respectivamente, los conjuntos de datos
residuo--cociente, el grupo abeliano de memoria de transporte, las
trayectorias registradas y las firmas de frontera.
Las realizaciones de esta firma se especifican con sus componentes en
las construcciones regional y terminal. Esta dependencia impide tratar
la firma como un parámetro libre una vez fijados los datos que la producen.

La carta elemental de residuo y cociente es explícita. Para
\(n\in\ZZ\), sean
\[
 r=1+((n-1)\bmod9),\qquad u=\frac{n-r}{9}.
\]
Entonces \(n=r+9u\), con \(r\in D_9\) y \(u\in\ZZ\).
La unicidad se sigue de que dos representantes de \(D_9\) congruentes
módulo \(9\) son iguales. Esta carta proporciona la reconstrucción
entera local. Las torres de precisión incorporan, además, las aplicaciones
de truncamiento y los cilindros desarrollados en
\ref{x_reciprocidad:sec:extension} y \ref{x_reciprocidad:sec:generacion}.

Sea \(\gamma:q\to q'\) una trayectoria observable. El transporte
de las fibras se expresa mediante aplicaciones
\[
 \mathsf H(\gamma):\mathsf H_q\to\mathsf H_{q'},\qquad
 \mathsf C(\gamma):\mathsf C_q\to\mathsf C_{q'},\qquad
 \mathsf{Hist}(\gamma):\mathsf{Hist}_q\to\mathsf{Hist}_{q'}.
\]
Las aplicaciones respetan identidades y composición, y
\(\mathsf C(\gamma)\) es un homomorfismo de grupos abelianos.
El incremento \(k_{\rm tr}(\gamma)\in\mathsf C_{q'}\) cumple
\begin{equation}
 \begin{split}
 k_{\rm tr}(\operatorname{id}_q)&=0,\\
 k_{\rm tr}(\gamma_2\circ\gamma_1)
 &=k_{\rm tr}(\gamma_2)+
   \mathsf C(\gamma_2)k_{\rm tr}(\gamma_1).
 \end{split}
 \label{x_reciprocidad:tpk-int:cociclo}
\end{equation}
El símbolo \(k_{\rm tr}\) designa aquí el cociclo de transporte;
la coordenada racional \(\kappa\) del registro \(K\) conserva
su notación en el capítulo de la clausura dodecafásica.

Las coordenadas transportables se actualizan por
\begin{equation}
 \begin{aligned}
 h'&=\mathsf H(\gamma)h,&
 c'&=\mathsf C(\gamma)c+k_{\rm tr}(\gamma),\\
 \lambda'&=\gamma\circ\lambda,&
 \sigma'&=\operatorname{Sig}_{q'}(h',c',\lambda').
 \end{aligned}
 \label{x_reciprocidad:tpk-int:accion-memoria}
\end{equation}
La segunda ecuación es una acción afín: el término lineal transporta
la memoria anterior y el cociclo añade el incremento producido por la
trayectoria. En la realización de cruce espacial de cada cursor por separado,
\(\mathsf C(\gamma)=I\) y \(k_{\rm tr}(\gamma)\) es la suma
de los vectores de \eqref{x_reciprocidad:tpk-int:borde}. Esta instancia concreta
vincula la ley abstracta con las operaciones cardinales ya definidas.

\begin{proposition}[Compatibilidad de la actualización con la concatenación]
\label{x_reciprocidad:tpk-int:composicion-memoria}
La actualización \eqref{x_reciprocidad:tpk-int:accion-memoria} por
\(\gamma_1\), seguida de la actualización por \(\gamma_2\),
coincide con la actualización por \(\gamma_2\circ\gamma_1\).
\end{proposition}
\begin{proof}
Para la coordenada afín, dos actualizaciones dan
\[
 c''=\mathsf C(\gamma_2)\mathsf C(\gamma_1)c
       +\mathsf C(\gamma_2)k_{\rm tr}(\gamma_1)
       +k_{\rm tr}(\gamma_2).
\]
La functorialidad de \(\mathsf C\) y
\eqref{x_reciprocidad:tpk-int:cociclo} convierten esta expresión en
\(\mathsf C(\gamma_2\circ\gamma_1)c+
k_{\rm tr}(\gamma_2\circ\gamma_1)\).
La afirmación para \(h\) se sigue de la composición de
\(\mathsf H\), y para \(\lambda\), de la asociatividad de las
trayectorias. La firma final coincide porque se evalúa sobre las mismas
tres coordenadas reconstruidas.
\end{proof}

El cilindro y su frontera completan esta ley mediante las relaciones
\(\operatorname{Ext}_\gamma\), definidas en
\ref{x_reciprocidad:tpk-int:bidireccionalidad}. Sus elementos determinan cuáles de
las actualizaciones anteriores conservan una prolongación admisible.
Los estados, junto con las ternas \((\gamma,x,x')\) admitidas por
esas relaciones, forman una categoría sobre las trayectorias observables:
\[
 (\gamma_2,x',x'')\circ(\gamma_1,x,x')
     =(\gamma_2\circ\gamma_1,x,x'').
\]
La identidad de \(x\) es \((\operatorname{id}_q,x,x)\).
La composición conserva a la vez el transporte aritmético y la
prolongabilidad de frontera. Esta es la estructura en la que se efectúan
los retornos con memoria y los refinamientos sucesivos.

\subsubsection{Registro dodecafásico, fase y memoria acumulada}
\label{x_reciprocidad:tpk-int:dodecafase}

Las \(108\) transiciones se distribuyen en las ventanas ordenadas
\[
 E_m=\{9m-8,\ldots,9m\},\qquad 1\leq m\leq12.
\]
El registro conserva el origen de fase, la orientación y los datos
producidos durante cada ventana. Su aplicación es
\[
 \mathcal R_{12}(x)=
 \bigl((E_m,\tau_m,\sigma_m,\kappa_m,\Lambda_m)\bigr)_{m=1}^{12},
\]
donde \(\tau_m\) es la subruta de la ventana, \(\sigma_m\) su
orientación, \(\kappa_m\) el incremento entero de frontera y
\(\Lambda_m\) el registro local de incidencias. Se adopta la carta
que se desarrollará en \ref{x_reciprocidad:subsec:k-registro-integral}; sus índices
de ventana corresponden a \(m=\beta+1\).
El dominio comprende las historias compatibles
del TPK y el codominio, sus registros temporales orientados. La
composición de segmentos y sus incrementos se rige por
\eqref{x_reciprocidad:tpk-int:accion-memoria}. El grupo \(C_{12}\) describe, en
cambio, la traslación del índice de ventana; \(\Theta_{108}\) conserva
el calendario de pasos. Esta distinción permite precisar qué se transforma
y qué se observa al completar un ciclo.

En las coordenadas \(\theta=9b+r\), con \(0\leq r<9\), el avance
de \(9\) pasos conserva \(r\) y aumenta \(b\) en una unidad.
La adición de fases introduce el cociente
\[
 c_0(r,r')=\left\lfloor\frac{r+r'}9\right\rfloor,
\]
cuya identidad de cociclo asegura la asociatividad de la composición.
En el calendario finito se conserva \([b]_{12}\); en un registro de
vueltas se conserva \(b\in\ZZ\), o \(b\in\mathbb N_0\) para
una historia de origen fijado. La sección~\ref{x_reciprocidad:sec:extension} demuestra
la sucesión exacta no escindida \eqref{x_reciprocidad:eq:extension108}. Su función es
describir la relación entre estas coordenadas, antes de emplearlas en la
construcción del estado firmado.

El registro firmado recibe los datos de fase y memoria mediante la
composición que se desarrolla en la sección~\ref{x_reciprocidad:sec:k-registro}:
\begin{equation}
 R_{\rm sgn}=\Pi_H\circ\mathcal E_{108}^{90,120}
                         \circ\mathcal R_{12}.
 \label{x_reciprocidad:tpk-int:registro-causal}
\end{equation}
Su dominio es el estado terminal con memoria; su salida pertenece a
\(\ZZ^{12}\). La transformación reversible entre ese estado firmado
y \(K\), así como la lectura racional de \(K\), se demuestran en
el capítulo correspondiente. Las operaciones de registro anteceden a
esas transformaciones. La reconstrucción de un registro desde sus
diferencias cíclicas verifica su recuperabilidad; la procedencia dinámica
debe conservarse mediante los eventos que lo produjeron.

Para la traza codificada \((d_t)\) conservada por el registro, el
lector de eventos cuenta los códigos \(\{2,4,6,8\}\) de cada
ventana con la orientación
\((+,+,+,-,-,-,+,+,+,-,-,-)\). Estos códigos y los símbolos
cardinales de \(\mathcal D\) son coordenadas distintas del
registro. El capítulo~\ref{x_reciprocidad:sec:k-registro} desarrolla este lector,
la descomposición integral \(1+5+6\), las congruencias que
caracterizan su imagen y su inversa exacta. La prueba de inversión
establece qué datos conserva esa carta; la definición de
\(\mathcal R_{12}\) especifica la historia orientada sobre la que
actúa.

\subsubsection{Incidencia de fases y construcción interna de los canales}
\label{x_reciprocidad:tpk-int:incidencia}

El selector determina tres subconjuntos de \(D_9\):
\[
 H_+=\{1,4,7\},\qquad H_-=\{2,5,8\},\qquad H_0=\{3,6,9\}.
\]
Sea \(H_{\rm act}=H_+\sqcup H_-\) el conjunto de fases activas y
sea \(O_{\rm ph}=D_9\setminus\{9\}\) la carta anterior al retorno
marcado. Denotemos por \(P_2(H_{\rm act})\) el conjunto de pares
no ordenados de fases activas. La incidencia interna está dada por
\begin{equation}
 \mathcal F^{\rm ph}_6=H_{\rm act}\times P_2(H_{\rm act}),\qquad
 \mathcal F^{\rm ph}_8=O_{\rm ph}\times P_2(H_{\rm act}),\qquad
 \Theta_{\rm ph}=D_9\times P_2(H_{\rm act}).
 \label{x_reciprocidad:tpk-int:fases90120}
\end{equation}
Los índices describen la cantidad de posiciones de la primera componente.
La construcción utiliza las fases y el origen del calendario ya fijados;
las realizaciones excepcionales posteriores transportarán estas
incidencias a sus propios dominios.

\begin{proposition}[Cardinales e incidencia orientada de fase]
\label{x_reciprocidad:tpk-int:cardinales}
Los conjuntos de \eqref{x_reciprocidad:tpk-int:fases90120} tienen cardinales
\(90,120,135\), respectivamente. Si
\[
 \partial_{\rm ph}=\Theta_{\rm ph}\times\{-1,+1\},\qquad
 E_{\rm ph}=\partial_{\rm ph}\sqcup\{\infty\},\qquad
 L_{\rm ph}=H_0\times P_2(H_{\rm act}),
\]
entonces
\[
 |\partial_{\rm ph}|=270,\quad |E_{\rm ph}|=271,\quad
 |L_{\rm ph}|=45,\quad
 |\partial_{\rm ph}\sqcup L_{\rm ph}|=315.
\]
Trasladar simultáneamente el origen de fase y el retorno marcado conserva
todos estos cardinales y sus inclusiones.
\end{proposition}
\begin{proof}
Hay seis fases activas, por lo que
\(|P_2(H_{\rm act})|=\binom62=15\). Los productos por
\(6,8,9\) dan \(90,120,135\). La orientación duplica \(135\);
la marca de retorno añade un elemento; las tres fases neutrales dan
\(3\cdot15=45\). Una traslación del calendario transporta cada
factor por una biyección y lleva el punto marcado a su imagen. Los
productos y las uniones disjuntas conservan así la incidencia.
\end{proof}

El cardinal \(271\) tiene aquí una realización por incidencia de
fase. La completación cúbica de la sección~\ref{x_reciprocidad:sec:extension} tiene
su propia construcción de la corona \(1000-729\). La correspondencia
entre ambas realizaciones debe conservar los mapas y las marcas que las
identifican, además del cardinal. Análogamente, el selector
\(\varepsilon_9\), la incidencia \((90,120)\) y el extractor
\(\mathcal E_{108}^{90,120}\) conservan los dominios y operaciones
que se acaban de distinguir.

\subsubsection{Inversión de trayectorias y prolongaciones conjugadas}
\label{x_reciprocidad:tpk-int:bidireccionalidad}

La bidireccionalidad comienza en una operación sobre trayectorias, antes
de asignarles valores escalares. Para
\(\gamma=(x_0;d_1,\ldots,d_n)\), con extremo \(x_n\), definimos
\[
 \operatorname{rev}\gamma
   =(x_n;\jmath d_n,\ldots,\jmath d_1).
\]
El origen de la trayectoria inversa es el extremo de la original.
La composición satisface
\[
 \operatorname{rev}^2\gamma=\gamma,\qquad
 \operatorname{rev}(\gamma_2\circ\gamma_1)
   =\operatorname{rev}\gamma_1\circ\operatorname{rev}\gamma_2,
\]
y el desplazamiento entero cambia de signo. Estas identidades se obtienen
invirtiendo el orden de los pasos y sustituyendo cada vector cardinal por
su opuesto. En un lazo se sigue
\(\operatorname{wind}(\operatorname{rev}\gamma)
=-\operatorname{wind}(\gamma)\).

Sobre los registros, la inversión requiere transportar asimismo el orden
de sus componentes y la orientación de sus eventos. La inversión de
ruta y una conjugación de firma son operaciones distinguibles: su
composición se especifica cuando se aplica a un registro concreto. En
particular, la involución regional que intercambia las componentes de
clausura y propagación conserva el eje de autoescala; su fórmula se
presenta en la sección~\ref{x_reciprocidad:mono-ampl:regiones}. Esta operación regional
actúa sobre otro dominio que la inversión de una trayectoria cardinal.

Las extensiones sobre una ruta \(r:q\to q'\) se expresan mediante
relaciones entre fibras de estados compatibles. Si \(\mathcal F_q\)
denota la fibra de estados completos sobre \(q\), se exige
\[
 \operatorname{Ext}_r\subseteq\mathcal F_q\times\mathcal F_{q'},
 \qquad \Delta_{\mathcal F_q}\subseteq
        \operatorname{Ext}_{\operatorname{id}_q},
\]
\[
 \operatorname{Ext}_{r_2}\circ\operatorname{Ext}_{r_1}
       \subseteq\operatorname{Ext}_{r_2\circ r_1}.
\]
Su composición conserva las evaluaciones de ambas hojas, la orientación,
los cocientes, el prefijo y la frontera. La inversión de la ruta no
implica por sí sola que toda relación de extensión sea una biyección.
La prolongación conjugada se define en el dominio donde el transporte de
esos datos satisface las compatibilidades correspondientes.

\subsubsection{Holonomía nonádica y continuidad del refinamiento}
\label{x_reciprocidad:tpk-int:holonomia}

Una vuelta de fase compone las transiciones sobre el estado con memoria:
\begin{equation}
 \Gamma_{9,k}=\mathcal U_{k+8}\circ\cdots\circ\mathcal U_k.
 \label{x_reciprocidad:tpk-int:gamma}
\end{equation}
Sobre una historia individualizada, esta composición funcional realiza
la relación de holonomía; en el dominio completo se conserva
\(\Gamma_{9,k}\subseteq\widehat X_k\times\widehat X_{k+9}\),
con las prolongaciones admisibles desarrolladas más adelante.
La base cíclica tiene nueve fases, mientras
la fibra conserva la trayectoria, los cocientes y la frontera. La
monodromía es la acción de una vuelta de esa base sobre las fibras; la
holonomía \(\Gamma_{9,k}\) expresa su realización por las transiciones
efectivas del TPK.

Si \(g\) observa la fase y \(M\) el número de vueltas conservadas,
\[
 g(\Gamma_{9,k}x)=g(x),\qquad M(\Gamma_{9,k}x)=M(x)+1.
\]
La primera identidad expresa clausura de fase y la segunda identifica
la nueva profundidad temporal. La prolongación coinductiva conserva
además sus cilindros y prefijos. Las restricciones entre profundidades
componen un sistema inverso; una sección compatible reúne sus
antecedentes sin reemplazarlos por la última observación. El desarrollo
de esta aritmética de historias y de su unidad ocupa
\ref{x_reciprocidad:hist:seccion}--\ref{x_reciprocidad:hist:doble-varianza}; la construcción de las
regiones y sus lectores se realiza en la sección~\ref{x_reciprocidad:sec:generacion}.

La transferencia de fase sobre emisiones ya construidas tiene otro
dominio. Sean \(\mathcal A_+\) y \(\mathcal A_\times\) los
alfabetos de firmas aditivas y multiplicativas, de cardinales \(18\)
y \(26\). Su producto identifica el catálogo \(\mathcal U_6\).
Para \(f:\mathcal A_+\times\mathcal A_\times\to\RR\),
las medias de hoja son
\[
 (E_+f)(s,e)=\frac1{18}\sum_{s'\in\mathcal A_+}f(s',e),\qquad
 (E_\times f)(s,e)=\frac1{26}\sum_{e'\in\mathcal A_\times}f(s,e').
\]
La construcción de estas emisiones y sus cardinales se efectúa mediante
la recurrencia de la sección~\ref{x_reciprocidad:sec:extension}. Una vez obtenidas,
el operador de transferencia se define sobre
\(\mathcal U_6\times C_9\) por
\[
 P_+=E_+,\quad P_-=E_\times,\quad P_0=I,\qquad
 \mathcal K_{\rm ph}((u,r),(v,r+1))=P_{M_{\rm ph}(r)}(u,v),
\]
con peso cero para las restantes transiciones de fase. El índice
\(r\) utiliza el representante positivo de la clase cíclica.
El orden de fases \((+1,-1,0)^3\) produce entonces la renovación
\[
 \mathcal K_{\rm ph}^{9}\big|_r=E_+E_\times.
\]
En efecto, ambas medias son idempotentes, conmutan y promedian factores
diferentes; su composición asigna peso \(1/(18\cdot26)=1/468\)
a cada emisión. Esta igualdad caracteriza un observable de transferencia
posterior. El transporte de las historias sigue dado por
\eqref{x_reciprocidad:tpk-int:gamma}, con su memoria y su frontera.

La jerarquía resultante establece la función de los capítulos siguientes.
La emisión finita construye el dominio y las regiones; las relaciones de
extensión conservan sus prefijos; la conexión nonádica los prolonga; el
registro dodecafásico determina los datos que intervienen en \(K\) y
en la clausura de \(\alpha\). La realización numérica y la realización
incidencial reciben esa estructura anterior. Este orden permite exponer
cada consecuencia en su capítulo conservando, desde el núcleo formal,
los objetos que la hacen posible.

% END OWNER



% BEGIN OWNER /Users/ruben/Documents/New project/output/EXTREMA_MEDIA_RAZON_RECIPROCIDAD_ES_EN_20260918/ARTICULO/ES/source/sections/memoria_resolvente.tex
\subsection{Memoria acumulada y resolventes del transporte}
\label{x_reciprocidad:subsec:memoria-resolvente}

El retorno del calendario no elimina los incrementos producidos durante
el recorrido. Para expresar esta diferencia operatoriamente, se construye
primero un registro acumulativo de eventos de la trayectoria y después su
serie generadora. La reducción de variables de memoria se realizará sobre
ese transporte, conservando también la fuente y el lector de salida.

\subsubsection{Eventos orientados y actualización del registro}

Sea \((d_t)\) la secuencia de códigos enteros de dirección conservada
por una trayectoria del TPK. El código \(d_t\) y la dirección cardinal
del cursor son coordenadas distintas; no se supone una identificación
entre sus alfabetos. Numeramos desde cero las ventanas
\[
 I_m=\{9m+1,\ldots,9m+9\},\qquad m\geq0.
\]
Durante el primer ciclo, \(I_m\) es la ventana \(E_{m+1}\) del registro
dodecafásico. Su orientación es
\(\sigma=(1,1,1,-1,-1,-1,1,1,1,-1,-1,-1)\).
Para una historia prolongada, \(\sigma_m\in\{-1,1\}\) es el signo
conservado en la ventana correspondiente. El evento firmado es
\begin{equation}
 a_m=\sigma_m\sum_{t\in I_m}
       \mathbf1_{\{2,4,6,8\}}(d_t),\qquad |a_m|\leq9.
 \label{x_reciprocidad:mem:evento}
\end{equation}
Cada sumando se incorpora cuando ocurre su transición. Esta definición
cuenta los eventos de la traza; no escoge direcciones mediante un valor
que deba obtenerse después.

Sean \(\mathbf e_0,\ldots,\mathbf e_{11}\) la base de \(\ZZ^{12}\)
y \(S\mathbf e_j=\mathbf e_{j+1\bmod12}\). Escribimos
\(R=S^{-1}\). El registro acumulado en coordenadas fijas satisface
\[
z_{m+1}=z_m+a_m\mathbf e_{m\bmod12}.
\]
Un prefijo sin historia acumulada empieza en \(z_0=0\); en otro
origen, \(z_0\) conserva el registro previo.
El marco que acompaña al calendario es \(y_m=S^{-m}z_m\).

\begin{proposition}[Actualización y retorno con memoria]
\label{x_reciprocidad:mem:prop-actualizacion}
El registro en el marco móvil cumple
\begin{equation}
 y_{m+1}=Ry_m+\eta_m,\qquad
 \eta_m=a_mR\mathbf e_0.
 \label{x_reciprocidad:mem:actualizacion}
\end{equation}
Para todo \(n\geq1\),
\begin{equation}
 y_n=R^ny_0+\sum_{j=0}^{n-1}R^{n-1-j}\eta_j.
 \label{x_reciprocidad:mem:iteracion}
\end{equation}
En particular,
\begin{equation}
 y_{12}=y_0+\sum_{j=0}^{11}a_j\mathbf e_j.
 \label{x_reciprocidad:mem:retorno12}
\end{equation}
\end{proposition}
\begin{proof}
Multiplicar la actualización de \(z_m\) por \(S^{-(m+1)}\) da
\(y_{m+1}=S^{-1}y_m+a_mS^{-1}\mathbf e_0\), porque
\(S^{-m}\mathbf e_{m\bmod12}=\mathbf e_0\).
La fórmula \eqref{x_reciprocidad:mem:iteracion} se obtiene por inducción.
Para \(n=12\), \(R^{12}=I\) y
\(R^{11-j}\eta_j=a_jR^{12-j}\mathbf e_0=a_j\mathbf e_j\),
lo que prueba \eqref{x_reciprocidad:mem:retorno12}.
\end{proof}

Así, doce ventanas restauran el marco y conservan el vector de eventos.
El recuento no permite recuperar por sí solo el orden de todos los pasos
dentro de cada ventana; ese orden permanece en la trayectoria registrada.

Para precisar el balance, definimos
\[
 D_3=I-S^3,\quad D_4=I-S^4,\quad
 B=D_3^{\mathsf T}D_3+D_4^{\mathsf T}D_4,
\]
\[
 \mathcal M(y)=\tfrac12y^{\mathsf T}By,
 \qquad q(y)=\sum_{j=0}^{11}y_j.
\]
La primera es una forma cuadrática del registro; la segunda registra
la suma de sus coordenadas. No se identifican aquí con energía ni carga físicas.

\begin{proposition}[Balance producido por el evento]
\label{x_reciprocidad:mem:prop-balance}
La actualización \eqref{x_reciprocidad:mem:actualizacion} satisface
\begin{equation}
 \begin{aligned}
 \mathcal M(y_{m+1})-\mathcal M(y_m)
   &=a_m(By_m)_0+2a_m^2,\\
 q(y_{m+1})-q(y_m)&=a_m.
 \end{aligned}
 \label{x_reciprocidad:mem:balance}
\end{equation}
\end{proposition}
\begin{proof}
Al desarrollar los productos,
\(B=4I-S^3-S^{-3}-S^4-S^{-4}\); por tanto
\(R^{\mathsf T}BR=B\) y \(B_{00}=4\).
Como \(y_{m+1}=R(y_m+a_m\mathbf e_0)\), la diferencia cuadrática
es \(a_m\mathbf e_0^{\mathsf T}By_m+\tfrac12a_m^2B_{00}\).
La suma de coordenadas es invariante bajo la permutación \(R\),
lo que da la segunda igualdad.
\end{proof}

\Needspace{20\baselineskip}
\subsubsection{Serie generadora y control uniforme de la memoria}

La serie conserva todos los incrementos, no únicamente su suma al final
de un ciclo. Con una variable formal \(u\), sean
\[
 Y(u)=\sum_{m\geq0}y_mu^m,\qquad
 H_\eta(u)=\sum_{m\geq0}\eta_mu^m.
\]

\begin{proposition}[Resolvente y cota de cola]
\label{x_reciprocidad:mem:prop-serie}
Se cumple la identidad formal
\begin{equation}
 Y(u)=(I-uR)^{-1}\bigl(y_0+uH_\eta(u)\bigr).
 \label{x_reciprocidad:mem:serie}
\end{equation}
Ambas series convergen para \(|u|<1\). Si \(\ell:\RR^{12}\to\RR\)
es lineal, \(L=\|\ell\|_{1\to\RR}\), \(M_0=\|y_0\|_1\) y
\(0<r<1\), entonces, para todo entero \(N\geq0\) y \(|u|\leq r\),
\begin{equation}
 \left|\sum_{m>N}\ell(y_m)u^m\right|
 \leq Lr^{N+1}\left(
 \frac{M_0+9(N+1)}{1-r}+\frac{9r}{(1-r)^2}\right).
 \label{x_reciprocidad:mem:cota}
\end{equation}
La serie de derivadas converge uniformemente en cada disco
\(|u|\leq r<1\).
\end{proposition}
\begin{proof}
Multiplicar \eqref{x_reciprocidad:mem:actualizacion} por \(u^{m+1}\) y sumar da
\(Y-y_0=uRY+uH_\eta\). El término constante de \(I-uR\) es
invertible, de modo que su inversa formal es \(\sum_{k\geq0}u^kR^k\).
Además, \(R\) preserva la norma \(\ell^1\) y
\(\|\eta_m\|_1\leq9\), de donde
\(\|y_m\|_1\leq M_0+9m\).
Para la cola se suman las mayorantes
\(L(M_0+9m)r^m\), usando
\[
 \sum_{m>N}r^m=\frac{r^{N+1}}{1-r},\qquad
 \sum_{m>N}mr^m=r^{N+1}
 \left(\frac{N+1}{1-r}+\frac r{(1-r)^2}\right).
\]
Por último, \(\sum_{m\geq1}Lm(M_0+9m)r^{m-1}<\infty\)
mayora la serie derivada y prueba su convergencia uniforme.
\end{proof}

La variable \(u\) cuenta ventanas de nueve transiciones. Su identificación
con la variable de otro lector requiere conservar ese reloj y el mapa
que relaciona ambas lecturas.

\subsubsection{Eliminación de memoria: operador, fuente y lector}

En un dominio donde la actualización está fijada como \(U:X\to X\),
el espacio vectorial libre \(V=\QQ^{(X)}\) sobre los estados admite el operador
\(T\delta_x=\delta_{U(x)}\). El reloj se incluye en \(x\) si la
transición depende de él. Esta extensión lineal conserva estados
distintos aunque compartan una fase observable.

Consideremos una descomposición \(V=V_0\oplus V_1\), donde \(V_1\)
es el sector auxiliar que se desea eliminar, y escribamos
\[
 T=\begin{pmatrix}A&B_1\\C&D\end{pmatrix},\qquad
 v=\binom{v_0}{v_1},\qquad \ell=(\ell_0,\ell_1).
\]
Aquí \(A:V_0\to V_0\), \(B_1:V_1\to V_0\),
\(C:V_0\to V_1\) y \(D:V_1\to V_1\).
Las identidades siguientes son formales; también valen donde los
resolventes convergen.

\begin{proposition}[Reducción de Schur con fuente y lector]
\label{x_reciprocidad:mem:prop-schur}
Sean
\[
 D_u=(I-uD)^{-1},\qquad
 \mathscr F(u)=I-uA-u^2B_1D_uC.
\]
La componente retenida de \((I-uT)^{-1}v\) es
\begin{equation}
 w_0=\mathscr F(u)^{-1}\bigl(v_0+uB_1D_uv_1\bigr),
 \qquad w_1=D_u(v_1+uCw_0).
 \label{x_reciprocidad:mem:schur-fuente}
\end{equation}
Su lectura completa satisface
\begin{equation}
 \begin{split}
 \ell(I-uT)^{-1}v
 ={}&(\ell_0+u\ell_1D_uC)
       \mathscr F(u)^{-1}(v_0+uB_1D_uv_1)\\
    &+\ell_1D_uv_1.
 \end{split}
 \label{x_reciprocidad:mem:schur-lector}
\end{equation}
\end{proposition}
\begin{proof}
Las dos ecuaciones de \((I-uT)w=v\) son
\(w_0=v_0+uAw_0+uB_1w_1\) y
\(w_1=v_1+uCw_0+uDw_1\).
Resolver la segunda y sustituirla en la primera da
\eqref{x_reciprocidad:mem:schur-fuente}; aplicar \(\ell_0\) y \(\ell_1\)
a ambas componentes da \eqref{x_reciprocidad:mem:schur-lector}.
Todas las inversas formales existen porque sus términos constantes
son identidades.
\end{proof}

El término \(u^2B_1D_uC=\sum_{k\geq0}u^{k+2}B_1D^kC\)
describe una salida hacia la memoria, \(k\) pasos dentro de ella y
un regreso. En cambio, \(u\ell_1D_uC\) lee la memoria antes de ese
regreso. Por ello, reducir sólo el operador no conserva necesariamente
la observación: también deben transformarse fuente y lector.

La diferencia es visible en el ciclo ya construido. Si \(P\) retiene
\(\mathbf e_0\), entonces \(PRP=0\), pero
\begin{equation}
 \left\langle\mathbf e_0,(I-uR)^{-1}\mathbf e_0\right\rangle
 =\frac1{1-u^{12}},\qquad
 \left\langle\mathbf e_{11},(I-uR)^{-1}\mathbf e_0\right\rangle
 =\frac u{1-u^{12}}.
 \label{x_reciprocidad:mem:ejemplo-ciclo}
\end{equation}
En efecto, \(R^n\mathbf e_0=\mathbf e_{-n\bmod12}\): las
dos series recogen, respectivamente, \(n\equiv0\) y
\(n\equiv1\pmod{12}\). El resolvente de la compresión es \(I\),
mientras la compresión del resolvente conserva los retornos omitidos.
Este ejemplo corresponde al registro de doce ventanas, no a la totalidad
del estado del TPK.

\subsubsection{Composición de segmentos y coherencia de la reducción}

Sean tres transportes afines composables
\(F_i(y)=R_i y+b_i\), con dominios y codominios compatibles.
La memoria producida por el recorrido completo es
\begin{equation}
 F_3F_2F_1(y)
 =R_3R_2R_1y+b_3+R_3b_2+R_3R_2b_1.
 \label{x_reciprocidad:mem:cociclo-tres}
\end{equation}
La prueba es la sustitución sucesiva de las tres fórmulas afines.
Agrupar primero los dos primeros segmentos o los dos últimos produce
la misma expresión; los factores transportan cada incremento desde
el punto donde se originó. En \eqref{x_reciprocidad:mem:actualizacion},
\(R_i=R\) y \(b_i\) es el evento de su ventana.

La eliminación de varias capas conserva una coherencia semejante.
Para hacerla explícita, descomponemos ahora
\(V=V_0\oplus V_1\oplus V_2\) y escribimos
\[
 T=\begin{pmatrix}
 A&B_1&B_2\\C_1&D_{11}&D_{12}\\C_2&D_{21}&D_{22}
 \end{pmatrix},\qquad Q_2=(I-uD_{22})^{-1}.
\]
Al eliminar la tercera componente quedan los cuatro bloques
\begin{equation}
 \begin{aligned}
 A'&=A+uB_2Q_2C_2,&
 B_1'&=B_1+uB_2Q_2D_{21},\\
 C_1'&=C_1+uD_{12}Q_2C_2,&
 D_{11}'&=D_{11}+uD_{12}Q_2D_{21}.
 \end{aligned}
 \label{x_reciprocidad:mem:schur-intermedio}
\end{equation}

\begin{proposition}[Transitividad de la eliminación]
Si \(D_*=(D_{ij})_{i,j=1}^2\), el complemento obtenido en dos etapas
coincide con el obtenido al eliminar juntas ambas memorias:
\begin{equation}
 \begin{split}
 I-uA'-u^2B_1'(I-uD_{11}')^{-1}C_1'
 ={}&I-uA\\
 &-u^2(B_1\ B_2)(I-uD_*)^{-1}\binom{C_1}{C_2}.
 \end{split}
 \label{x_reciprocidad:mem:schur-transitivo}
\end{equation}
\end{proposition}
\begin{proof}
Resolver la tercera ecuación de \((I-uT)w=v\) y sustituirla en las
dos primeras produce \eqref{x_reciprocidad:mem:schur-intermedio}. Eliminar después la
segunda componente da el miembro izquierdo de
\eqref{x_reciprocidad:mem:schur-transitivo}. Resolver a la vez las dos ecuaciones de
memoria da el derecho, por la proposición~\ref{x_reciprocidad:mem:prop-schur}.
La solución formal es única porque \(I-uT\) tiene término constante
\(I\). En ambos procedimientos se transforman también la fuente y
el lector según \eqref{x_reciprocidad:mem:schur-fuente}--\eqref{x_reciprocidad:mem:schur-lector}.
\end{proof}

El balance \eqref{x_reciprocidad:mem:balance} describe los incrementos del registro.
Las identidades anteriores permiten transportarlos al reutilizar la memoria
en la clausura dodecafásica; no determinan por sí solas coeficientes
analíticos adicionales.

% END OWNER

